% !PRINTED 114
\weberpage{114}

Accordingly we introduce the three functions:
\weq{9}{
\begin{aligned}
f(\omega) &= \frac{e^{-\frac{\pi i}{24}}\eta\left(\dfrac{\omega+1}{2}\right)}{\eta(\omega)},\\
f_1(\omega) &= \frac{\eta\left(\dfrac{\omega}{2}\right)}{\eta(\omega)},\\
f_2(\omega) &= \sqrt{2}\,\frac{\eta(2\omega)}{\eta(\omega)}.
\end{aligned}
}

Then by (6), (7), (8)
\weq{10}{
\begin{aligned}
\vartheta_{00} &= \eta(\omega)f(\omega)^2,\\
\vartheta_{01} &= \eta(\omega)f_1(\omega)^2,\\
\vartheta_{10} &= \eta(\omega)f_2(\omega)^2,
\end{aligned}
}
and from (8) and (9) it follows that $f(\omega)$, $f_1(\omega)$,
$f_2(\omega)$ are real and positive for a purely imaginary $\omega$.
(10) agrees with § 24, (10), and the functions $f$, $f_1$, $f_2$
introduced there are the same as these. But from (10) the relations
follow easily [(2) and § 21 (14)]:
\weq{11}{
\begin{aligned}
f(\omega)^8 &= f_1(\omega)^8 + f_2(\omega)^8,\\
\sqrt{2} &= f(\omega)f_1(\omega)f_2(\omega).
\end{aligned}
}

With the help of these formulae one can express each of the three
functions $f(\omega)^8$, $f_1(\omega)^8$, $f_2(\omega)^8$ by square
roots in terms of any other. The formula flowing from (11) leads to
this:
\weq{12}{
\begin{aligned}
{[f_1(\omega)^8 - f_2(\omega)^8]^2} &= f(\omega)^{16} - \frac{64}{f(\omega)^8},\\
{[f(\omega)^8+f_2(\omega)^8]^2} &= f_1(\omega)^{16} + \frac{64}{f_1(\omega)^8},\\
{[f(\omega)^8+f_1(\omega)^8]^2} &= f_2(\omega)^{16} + \frac{64}{f_2(\omega)^8}.
\end{aligned}
}

For the linear fundamental transformations of the functions $f$ one
obtains first of all from (4) and (9):
\weq{13}{
\begin{aligned}
f(\omega+1) &= e^{-\frac{\pi i}{24}}f_1(\omega),\\
f_1(\omega+1) &= e^{-\frac{\pi i}{24}}f(\omega),\\
f_2(\omega+1) &= e^{\frac{\pi i}{12}}f_2(\omega).
\end{aligned}
}
